\documentclass{article}

% if you need to pass options to natbib, use, e.g.:
%     \PassOptionsToPackage{numbers, compress}{natbib}
% before loading neurips_2026

\usepackage[final]{neurips_2026}
\makeatletter
\providecommand{\@trackname}{}
\makeatother
\usepackage{graphicx}

\usepackage[utf8]{inputenc} % allow utf-8 input
\usepackage[T1]{fontenc}    % use 8-bit T1 fonts
\usepackage{hyperref}       % hyperlinks
\usepackage{url}            % simple URL typesetting
\usepackage{booktabs}       % professional-quality tables
\usepackage{amsfonts}       % blackboard math symbols
\usepackage{amsmath}        % math environments
\usepackage{nicefrac}       % compact symbols for 1/2, etc.
\usepackage{microtype}      % microtypography
\usepackage{xcolor}         % colors

\title{The NDC Data Explorer: A Design Case Study of Citation-Grounded AI for Climate Policy Briefing at the UNFCCC}

\author{%
  Tseten Sherpa$^{1,2}$ \\
  $^{1}$ School for Environment and Sustainability, University of Michigan \\
  $^{2}$ Mitigation Data Systems, UNFCCC, Bonn, Germany \\
  Correspondence to: tsherp@umich.edu
}

\begin{document}

\maketitle

\begin{abstract}
Governments submit Nationally Determined Contributions (NDCs) to the UNFCCC every five years, documenting climate mitigation and adaptation commitments under the Paris Agreement. These submissions are long, heterogeneous, and produced under significant capacity constraints, particularly in developing-country delegations. This paper is a \emph{design case study} of the NDC Data Explorer, a decision-support cockpit built during a UNFCCC internship so that secretariat staff and government negotiators can place Uganda's Updated NDC (September 2022) next to independently measured emissions and query that view in natural language. The system is a React and Express application that pulls live Climate TRACE inventories, encodes official BAU-relative sector ceilings, and exposes an assistant that is allowed to quote only numbers from a typed fact ledger built from the current dashboard. Complementary modules---a MapLibre three-dimensional source map of Uganda, a Climate Policy Radar document corpus, and an explicitly labelled indicative climate-finance screen---were piloted with Uganda as the sole fully wired testbed. A June 2026 demonstration and subsequent recognition from the Executive Secretary are recorded as institutional reception. They are not measurements of accuracy, time saved, or decision quality. We do not claim to have shown that AI applications help governments with climate and energy policy. We document one scoped architecture and state four design principles as \emph{hypotheses to be tested}: that an assistant helps only when limited to retrieval, comparison, and summarization over verified sources; when live observations are never silently blended with pledges or modelled figures; when numeric claims are checkable and near-miss fabrications are not certified; and when progress against a BAU-relative ceiling is scored against the no-policy path in the observation year, not against the 2030 BAU point.
\end{abstract}

\section{Introduction}
\label{sec:intro}

Under Article~4 of the Paris Agreement, every Party to the UNFCCC submits and periodically updates a Nationally Determined Contribution describing its emissions-reduction targets and the policies intended to achieve them [5]. NDCs are simultaneously technical, legal, and political documents. Tracking their content, comparing commitments across countries, and reconciling stated targets against independently measured emissions is a labor-intensive task that falls on a small number of secretariat staff and national negotiators, many of whom work with limited technical infrastructure. Recent reviews of natural language processing in this space note that climate policy documents are long, information-dense, inconsistently structured, and difficult to navigate without substantial manual effort and domain expertise [3].

The operational problem is sharper than document length. Ambition is already written down; what ministries often lack is a live picture of delivery. Official NDC PDFs, national greenhouse-gas inventories, satellite-based third-party estimates, climate-finance tables, and sector strategies live in separate files, on different update cycles, and in different accounting frames. A negotiator who needs to know whether transport is on a path consistent with a 2030 ceiling, where the largest observed sources sit, and what the NDC actually pledged must currently assemble that answer by hand.

This capacity gap is the setting in which AI applications are proposed as a remedy, and the setting in which their risks are most consequential. A government relying on a hallucinated summary of its own emissions commitments, or of another Party's, is a diplomatic liability. This paper describes a system, the NDC Data Explorer, built as a prototype with which to \emph{pose}---not yet answer---the question of whether a narrowly scoped assistant can be useful in that work while remaining accountable to its evidence.

We make three design contributions, not empirical findings. First, we describe a deployed Uganda testbed that compares official BAU-relative NDC ceilings to live Climate TRACE inventories [7] without forcing the two accounting systems to agree, and we correct an optimistic progress definition that compared today's emissions to the 2030 no-policy point. Second, we specify a fact-ledger pattern for numeric policy question-answering and a citation-resolver policy that refuses to treat a near-miss numeral as verified. Third, we record four design principles as hypotheses for later measurement rather than as conditions the pilot has established.

\section{Background: AI in the climate policy pipeline}
\label{sec:background}

A growing body of work has applied large language models (LLMs) to climate and sustainability policy text, including classification of NDC content against UNFCCC taxonomies, summarization of climate action plans, and synthesis of scientific literature for policymakers. An empirically grounded assessment of LLMs on climate and sustainability policy documents found that these tools are useful for processing, classifying, and summarizing heterogeneous text, but identified persistent limits around interpretability, external validity, replicability, and usability that can threaten decision-making if left unaddressed [2]. A broader narrative review of LLM use in climate-focused social research similarly finds real opportunities for expanding the scale of climate research and communication, alongside risks that, if unmanaged, could turn the technology into another obstacle to climate action rather than an accelerant [3]. Separately, generative agent-based approaches have been proposed to help policymakers anticipate public acceptability of proposed climate policies, using LLMs as approximate models of public response [4].

A recurring theme is that the value of LLMs in policy settings depends on whether outputs are grounded in verifiable source material. Retrieval-augmented generation (RAG) conditions a language model's output on passages retrieved from a fixed corpus rather than relying purely on parametric memory [1], and has become a standard mitigation for hallucination in knowledge-intensive applications. The NDC Data Explorer is related to this line of work but does not implement classical vector RAG for its primary assistant. Dashboard answers are conditioned on a structured, typed ledger of numbers that already appear on screen, each paired with a resolvable source URL. Document-level assistance, by contrast, extracts text from an allowlisted PDF and asks the model to cite page markers. Both designs share RAG's motivation---do not let fluency substitute for evidence---while choosing a grounding mechanism matched to the claim type: numbers from a live dashboard versus passages from a named policy file.

Independently of language models, open emissions inventories now make it possible to place national pledges next to third-party observations at sector and source resolution. Climate TRACE synthesizes satellite observation, activity data, and sector models into a regularly updated global inventory [7, 8]. That independence is useful for accountability and planning; it is not a substitute for a Party's official greenhouse-gas inventory. The two series can differ in sector definitions, gas coverage, treatment of land use, and update lag. A responsible interface must show the comparison and the discrepancy, not average them away. Climate Policy Radar's open corpus of climate laws, policies, and international submissions [9] likewise makes it possible to keep the NDC and related national documents in the same workspace as the numbers, as a curated export rather than an implicit web search.

This combination---structured pledges, an independent inventory, and a document snapshot---defines a narrower AI problem than ``apply an LLM to climate policy.'' The design question is whether a model can be prevented from mixing those layers. Classical RAG retrieves text; it does not, by itself, stop a model from adding a plausible MtCO$_2$e figure that appeared in pretraining. A typed ledger is a more pedantic tool. Whether that pedantry is sufficient is an empirical question this paper does not answer.

\section{Setting: Uganda as a testbed}
\label{sec:setting}

The explorer was piloted on Uganda's Updated NDC of September 2022 [6], chosen in coordination with UNFCCC secretariat staff. Uganda is a least-developed country whose mitigation ambition is partly conditional on international finance and support. The 2022 update is BAU-relative: emissions are expected to grow with the economy, and sector targets are expressed as a percentage below a 2030 no-policy baseline rather than as an absolute cut from 2015. Misreading that framing is a common source of dashboard error. A 2030 ceiling can sit \emph{above} the 2015 inventory; ``progress'' is then position on a BAU path, not a reduction from the base year.

Economy-wide, Uganda pledges a 24.7\% reduction below 2030 BAU, split into 5.9\% unconditional and 18.8\% conditional support, corresponding to 112.1~MtCO$_2$e in 2030. Sector ceilings used by the cockpit are taken from the same submission and encoded as versioned configuration (Table~\ref{tab:ndc-targets}).

\paragraph{The 2030-BAU comparator is optimistic.}
The June 2026 prototype scored BAU-cap emissions targets with a form common in tracking tools [10]:
\begin{equation}
\label{eq:progress-2030}
  \mathrm{progress}_{2030}(\%)
  \;=\;
  100 \times
  \frac{\mathrm{BAU}_{2030} - y_{t}}{\mathrm{BAU}_{2030} - C_{2030}},
\end{equation}
clamped to $[0,100]$, where $y_{t}$ is the latest observed value and $C_{2030}$ is the NDC ceiling. Because $\mathrm{BAU}_{2030}$ is a no-policy projection for a future year, it lies above today's inventory for every growing sector. Substituting $y_{t}$ into \eqref{eq:progress-2030} therefore credits a country with a large share of the 2030 wedge even if emissions have followed BAU to date. Inaction looks like partial delivery.

\paragraph{Contemporaneous BAU.}
The quantity that should be compared with $y_{t}$ is the no-policy level in year $t$, not in 2030. The NDC does not publish a full BAU timeseries, so we take a transparent linear interpolant from the base-year inventory $B_{t_0}$ to the 2030 BAU point, and the same interpolant for the ceiling:
\begin{align}
\label{eq:bau-t}
  \mathrm{BAU}_{t}
  &=
  B_{t_0} + (\mathrm{BAU}_{2030}-B_{t_0})\,\frac{t-t_0}{2030-t_0}, \\
\label{eq:cap-t}
  C_{t}
  &=
  B_{t_0} + (C_{2030}-B_{t_0})\,\frac{t-t_0}{2030-t_0}.
\end{align}
At $t=t_0$ the denominator of the progress ratio is zero; we define progress as undefined (displayed as unknown) in the base year. For $t>t_0$,
\begin{equation}
\label{eq:progress-t}
  \mathrm{progress}_{t}(\%)
  \;=\;
  100 \times
  \frac{\mathrm{BAU}_{t} - y_{t}}{\mathrm{BAU}_{t} - C_{t}},
\end{equation}
again clamped to $[0,100]$. Observed values above $\mathrm{BAU}_{t}$ score as off-track relative to the contemporaneous no-policy path; values at or below $C_{t}$ score as complete on this metric. Linear interpolation is itself an assumption, stated as such: it is a tracking convention, not a reconstruction of Uganda's official BAU model. True-reduction targets (2030 below the base-year inventory) continue to use the conventional gap-from-baseline form. Agriculture climate-smart-agriculture adoption, forest cover, wetlands coverage, electricity capacity, and electricity access are tracked as national indicators, not as Climate TRACE sectors.

A transport example shows both the BAU-relative framing and the 2030 bias. The NDC records a 2015 transport inventory of 4.2~MtCO$_2$e, a 2030 BAU of 9.6, and a 2030 ceiling of 6.8 (29\% below BAU). The ceiling is \emph{higher} than the base year, so a dashboard that treated $4.2\to 6.8$ as a failed cut would report that Uganda had already missed a target it had not set. That part of the design is correct. The error is the year of the BAU. Suppose the latest observation is 7.5~MtCO$_2$e in 2024. Equation~\eqref{eq:progress-2030} yields $100\times(9.6-7.5)/(9.6-6.8)\approx 75\%$: most of the way to the goal after doing little more than remaining below a distant 2030 BAU. Equations~\eqref{eq:bau-t}--\eqref{eq:progress-t} give $\mathrm{BAU}_{2024}=7.44$ and $C_{2024}=5.76$, so progress is $100\times(7.44-7.5)/(7.44-5.76)<0$, clamped to $0\%$: the observation sits on (slightly above) the interpolated no-policy path. The same contrast applies to AFOLU (77.6 baseline, 122.2 BAU, 91.8 cap) and stationary energy. Waste is the near-flat case (2.08 to a 2.09 cap against 3.19 BAU): being near the 2015 level can be consistent with the pledge once the comparator is $\mathrm{BAU}_{t}$, not evidence of inaction and not automatic credit against 2030.

The prototype that was demonstrated still implemented \eqref{eq:progress-2030}. Equation~\eqref{eq:progress-t} is the corrected definition this case study now recommends; replacing the shipped metric is a specified change, not a result already measured in use.

Geography is national or district. National Climate TRACE series begin in 2015; district series begin in 2021. Fifty-six Uganda districts are available via Climate TRACE GADM identifiers (national \texttt{UGA}; for example Kampala \texttt{UGA.16\_1}). NDC targets are national. District numbers are therefore shown as observed context only: the interface does not assign a district a pass/fail score. A Kampala timeseries can inform siting and briefing; it cannot be read as ``Kampala is off track for the NDC,'' because no such district target exists. That is a design choice, not an unfinished feature. Other countries appear in a country gate as scaffolding; they are not wired to live data.

\begin{table}[t]
  \caption{Uganda Updated NDC (September 2022) mitigation ceilings used by the cockpit. Sector targets are BAU-relative. Agriculture emissions are a Climate TRACE-tracked component of AFOLU and have no standalone NDC 2022 ceiling. Units MtCO$_2$e except where noted.}
  \label{tab:ndc-targets}
  \centering
  \small
  \begin{tabular}{llrrrl}
    \toprule
    ID & Sector & 2015 base & 2030 BAU & 2030 cap & CT \\
    \midrule
    t0 & Economy-wide & --- & --- & 112.1 & --- \\
    t1 & AFOLU ($-24.9\%$ BAU) & 77.6 & 122.2 & 91.8 & yes \\
    t4 & Energy, stationary ($-18.8\%$) & 5.66 & 12.44 & 10.10 & yes \\
    t5 & Transport ($-29\%$) & 4.2 & 9.6 & 6.8 & yes \\
    t6 & Waste ($-34.8\%$) & 2.08 & 3.19 & 2.09 & yes \\
    t7 & IPPU ($-14\%$) & 0.57 & 1.0 & 0.86 & yes \\
    t8 & Agriculture (CSA indicator) & --- & --- & --- & ind. \\
    \bottomrule
  \end{tabular}
\end{table}

\section{System overview}
\label{sec:system}

The NDC Data Explorer is a web-based decision-support cockpit, not a chatbot wrapped around a search index and not official MRV submission software. A React and TypeScript single-page application (Vite, Tailwind, MapLibre GL) talks to an Express API under \texttt{/api/v1}. Production is deployed on Vercel. Live emissions and progress come from Climate TRACE API~v7 (gas field \texttt{co2e\_100yr}); the only numeric transform applied for display is tonnes to MtCO$_2$e. NDC target definitions, sector-to-slug mapping, and the mitigation catalogue are versioned in the repository. Optional Postgres stores ingested indicator observations; it does not replace Climate TRACE MtCO$_2$e series. Browser storage holds personal activity drafts and a role preference. Role labels (ministry officer, MRV officer, senior decision-maker, and others) shape navigation defaults; they are not remote authentication.

Figure~\ref{fig:architecture} summarizes the data path. Three upstream classes feed the API: live Climate TRACE inventories, bundled Uganda NDC 2022 configuration, and a Climate Policy Radar export of Uganda documents and multilateral climate-fund projects. The cockpit assembles those into a dashboard, map, document library, and finance screen. When a user asks a question, the client builds a fact ledger from the current view and sends it with the query to an OpenAI chat-completions model; the server returns cited prose and a confidence flag.

\begin{table}[t]
  \caption{Architecture of the NDC Data Explorer. The cockpit is the system of record for what the assistant may say: the model never browses the open web.}
  \label{fig:architecture}
  \centering
  \small
  \begin{tabular}{ccc}
    \toprule
    Climate TRACE API v7 & Uganda NDC 2022 config & CPR / MCF document export \\
    \midrule
    \multicolumn{3}{c}{$\downarrow$ \quad Express API (emissions, documents, ingest) \quad $\downarrow$} \\
    \multicolumn{3}{c}{$\downarrow$ \quad React cockpit (dashboard, map, finance, documents) \quad $\downarrow$} \\
    \multicolumn{3}{c}{fact ledger $\rightarrow$ OpenAI chat (quote-only) $\rightarrow$ cited prose} \\
    \bottomrule
  \end{tabular}
\end{table}

A cross-cutting honesty model, summarized in Table~\ref{tab:honesty}, is as important as the architecture. Every user-visible number is labelled live, bundled-from-NDC, snapshot, indicative, or local-only. Indicative modules may inform screening; they are not permitted to overwrite Climate TRACE totals or NDC ceilings.

\begin{table}[t]
  \caption{Data-honesty taxonomy. The cockpit refuses to treat these classes as interchangeable.}
  \label{tab:honesty}
  \centering
  \small
  \begin{tabular}{p{0.18\columnwidth}p{0.34\columnwidth}p{0.36\columnwidth}}
    \toprule
    Class & What users see & Provenance \\
    \midrule
    Live & Sector and district MtCO$_2$e, timeseries, map sources & Climate TRACE API v7; tonnes to MtCO$_2$e only \\
    Bundled official & 2030 ceilings, BAU, conditionality, indicator targets & Uganda Updated NDC (Sept.\ 2022) \\
    Snapshot & $\sim$207 policy documents; $\sim$167 MCF projects & Climate Policy Radar export, not a live API \\
    Indicative & MACC / fund fit; KCI analogies; 2030 trend; risk & Catalogue assumptions or seed models \\
    Local-only & User activities; role preference & Browser storage; not a shared ledger \\
    \bottomrule
  \end{tabular}
\end{table}

\section{Methods}
\label{sec:methods}

\subsection{Emissions dashboard}
\label{sec:dashboard}

The primary analysis surface is a five-column NDC cockpit: official targets, observed Climate TRACE series, progress toward the selected ceiling, mitigation options, and NDC-traceable activities. Users switch sector (AFOLU, energy, transport, IPPU, agriculture, waste) and geography. Climate TRACE subsector slugs are aggregated through a fixed sector map; a reconciliation invariant requires that mapped sectors plus a small unmapped remainder (for example mineral extraction) explain the national total. An on-screen accuracy strip reports API health and the reconciliation residual. Named top emitting sources are spatially located assets and are \emph{not} required to sum to the sector total: spatially uncertain mass appears in totals only.

The dashboard does not invent missing inventory years and does not interpolate Climate TRACE gaps. When Climate TRACE and the NDC inventory-scale baseline diverge---as they systematically can, because one is a satellite-and-model inventory and the other is a national GHG inventory---the interface shows a warning rather than coercing equality. A planning tool that hid that discrepancy would train users to treat third-party estimates as official MRV. Exports (Excel, a plain-language PDF briefing, and a CRT/BTR-style CSV) inherit the current geography label so a district view cannot be mistaken for a national submission table.

Numeric integrity is checked rather than assumed. Unit tests lock the sector-slug map, the tonnes-to-MtCO$_2$e conversion, and a no-mock-fallback rule for production routes. A live verification script compares dashboard aggregates to Climate TRACE v7 for a reference year and flags residuals above a small MtCO$_2$e tolerance. Those checks do not make the inventory ``correct''; they make the application's transforms auditable. They also do not yet lock \eqref{eq:progress-t}: the shipped tests encode \eqref{eq:progress-2030}. Updating that suite is part of the progress correction.

\subsection{Spatial layer}
\label{sec:map}

The emissions map is a MapLibre GL three-dimensional satellite and terrain view of Uganda, not a global globe. Geolocated Climate TRACE sources are drawn as bubbles sized with the square root of MtCO$_2$e, coloured by sector, for inventory years 2021--2025. The renderer caps displayed points (on the order of four thousand) even if the API returns more; users are not shown a silently truncated ``complete'' inventory. A district choropleth provides administrative context. The map answers \emph{where observed emissions sit}; pass/fail against the NDC remains a national dashboard question.

\subsection{Fact-ledger NDC AI}
\label{sec:ndc-ai}

The core natural-language interface is an OpenAI chat-completions assistant (default model configurable; deployment used a current chat model) that answers questions about the \emph{current dashboard state}. It is not an embedding retriever over a document store and not a web-enabled agent.

Before each call, the client constructs a fact ledger: a list of typed records with an identifier, value, claim, source URL, viewer URL, and domain. Identifiers distinguish inventories from pledges (for example \texttt{fact\_trace\_transport\_latest\_2024} versus \texttt{fact\_ndc\_t5\_target}). Typical rows include, for each dashboard sector, the latest Climate TRACE value and year, the last five observed years, the NDC 2030 ceiling and BAU, derived progress percent, the national country total, and Uganda's global rank. Selected-target rows add the official baseline, target, and a plain-language gloss of the pledge. Domains are Climate TRACE or UNFCCC. Source URLs point at the Climate TRACE v7 emissions or rankings endpoints that produced the number, or at the UNFCCC NDC Registry PDF for Uganda [6]. The public Climate TRACE inventory page is attached as a human viewer URL, not as a second numeric source.

The system prompt requires the model to (i)~quote only values that appear in the quotable-facts list, (ii)~attach the corresponding fact identifiers to each paragraph, (iii)~keep Climate TRACE and NDC numbers on separate citation lists, and (iv)~say that a figure is unavailable rather than estimate it. Output is constrained to a JSON schema: title, confidence, headed sections of short prose paragraphs, a one-sentence disclaimer, and follow-up questions. Free-text questions are capped at 500 characters. Four canned actions---progress check, gap analysis, sector emissions, and priorities---use the same ledger with a narrower task line, so a briefing officer can obtain a structured note without composing a prompt.

Consider the question ``Are we on track for transport?'' The ledger already contains the NDC transport ceiling, BAU, baseline, the latest TRACE estimate, and the derived progress percent. A valid answer may compare those figures and cite NDC facts for the pledge and TRACE facts for the observation. It may not invent a 2015--2030 reduction rate, a district pass/fail, or a finance need, because those numbers are not in the ledger. That is the intended scope: the assistant is allowed to narrate the cockpit, not to enlarge it. Whether users treat it that way has not been measured.

\paragraph{Near-miss certification is a failure mode.}
After generation, a deterministic resolver maps cited identifiers to URLs. The shipped resolver also tried numeric matching: if a paragraph number lay within an absolute tolerance of $0.06$ or a relative tolerance of $2.5\%$ of a ledger value, the checker treated the claim as that fact and attached the real source URL. Keyword inference was a further fallback when no identifier and no numeric match existed. The first of those rules certifies a fabrication. A model that invents $7.4$ when the ledger holds $7.5$ (about $1.3\%$ away) receives a Climate TRACE or UNFCCC link and looks verified.

The direction of the fix is the reverse of matching the model's number to the ledger. A near-match must not mint a citation. The intended policy is: (i)~if the model cited a fact identifier, display the \emph{exact} ledger value in place of the generated numeral when they differ, so the user sees the verified figure; (ii)~if there is no identifier, do not search for a nearby ledger value in order to bless the sentence---leave the model's number, attach no source from that match, and mark the paragraph unverified; (iii)~any paragraph that still contains a number absent from the ledger after exact comparison (not $2.5\%$ comparison) forces the answer's confidence to low (Table~\ref{fig:ai-pipeline}). Keyword inference remains a residual risk if it attaches a URL to an unrelated sentence; it should not be used to recover a missing identifier when numbers fail an exact check.

This paper does not report how often the shipped checker would pass a close-but-wrong number. That measurement requires an adversarial suite: queries designed to elicit rounded, transposed, or interpolated MtCO$_2$e figures, scored against a labelled ledger. Until those failure rates exist, citation grounding is a design hypothesis, not a validated property.

\begin{table}[t]
  \caption{NDC AI pipeline. The model is not the system of record. Near-miss numerals are overwritten or flagged; they are not certified with a real URL.}
  \label{fig:ai-pipeline}
  \centering
  \small
  \begin{tabular}{ccccc}
    \toprule
    Dashboard state & $\rightarrow$ & Typed fact ledger & $\rightarrow$ & LLM (quote-only) \\
     & & & $\rightarrow$ & Exact citation resolver \\
     & & & $\rightarrow$ & Cited answer or low-confidence flag \\
    \bottomrule
  \end{tabular}
\end{table}

\subsection{Policy documents}
\label{sec:documents}

A second assistant operates over documents rather than dashboard numbers. The library is a Climate Policy Radar export: on the order of 207 Uganda rows (laws, executive plans, UN submissions, fund projects), a passage index over key texts (NDC, NDPIV, NBSAP, agroforestry and climate regulations), and a multilateral climate-fund (MCF) project table of on the order of 167 Uganda-linked records. This is a snapshot rebuilt from CSV, not a live CPR API.

Document analysis fetches the PDF over HTTPS, extracts text with page markers, truncates to 8{,}000 characters, and asks the model for an executive summary, key items, targets, or recommendations with page citations of the form $[p.N]$. Truncation is a limitation: a long NDC annex will not be fully in context. We prefer a partial, page-addressable reading to an unbounded fetch that pretends to have read the whole file. Server-side fetch is host-allowlisted to the Climate Policy Radar CDN; redirects are disabled. The allowlist is a trust boundary: the model may read only files the secretariat already treated as in-corpus. Passage search over the key-document index is a separate, non-generative path: keyword or topic filters return grouped passages without calling the language model, which is the right tool when the user already knows which instrument they mean.

\subsection{Climate finance}
\label{sec:finance}

Mitigation ambition in many NDCs, including Uganda's, is explicitly conditional on international support. The finance module screens catalogue options by indicative abatement and cost: default assumptions are a carbon price of \$20/t, a ten-year lifetime, and a 10\% discount rate, which users can move. Outputs include a marginal abatement cost curve, annualized cost, and a fund-window match (GCF readiness, simplified approval, and full proposals; GEF; World Bank IDA and related windows) given Uganda's LDC status. Sequencing advice is deliberately pedestrian: pick a window whose scale matches the measure, then write toward that window, rather than retrofitting a finished concept note. MCF project records can be searched from the same screen. The UI is forbidden from implying audited project costs or investment advice. Live sector gaps from Climate TRACE may appear as context; they do not become cost figures.

\subsection{Adjacent capabilities}
\label{sec:adjacent}

Three further modules are shipped and labelled so they cannot be mistaken for the live inventory. An AI 2030 view fits a lightweight trend on Climate TRACE history; it is a planning sketch, not a forecast of record. A policy-impact wizard matches interventions to UNFCCC Katowice Committee of Experts on the Impacts of the Implementation of Response Measures (KCI) case analogies and Transition Element Framework labels; matches are indicative socio-economic analogies, not Uganda-specific attribution. Mapped ingest lets an MRV officer attach observations to indicator targets (forest cover, electricity, climate-smart agriculture, wetlands) in Postgres; those observations never replace Climate TRACE MtCO$_2$e. Climate-risk layers remain illustrative seed data.

\section{Deployment}
\label{sec:deployment}

Development included supervised collaboration with the UNFCCC's Data Science Manager, as well as external input connected to Climate TRACE. The internship is based at the Secretariat in Bonn and is being carried out remotely from the United States through November 2026. The application is hosted on Vercel; a live prototype is available at \url{https://ndc-data-explorer-e051f914.vercel.app/select-country}. The prototype was demonstrated at the UNFCCC's June 2026 climate meetings along a live path of dashboard, map, documents, finance, and NDC AI, and received direct recognition from Executive Secretary Simon Stiell. That reception is a fact about the internship; it is not a substitute for user-study outcomes, query-level accuracy, or time-on-task. The project is being evaluated for longer-term development and funding beyond its original internship scope.

The demonstration framing used with stakeholders was deliberately modest: governments do not lack written ambition; they lack a live delivery picture. Features most emphasized in those conversations were the accuracy strip (Climate TRACE health and reconciliation), the refusal to force TRACE and inventory numbers into equality, and citation grounding on the assistant. Country coverage remains Uganda-complete, with a gate that can later attach other Parties without changing the honesty model.

\section{Discussion: Design principles as hypotheses}
\label{sec:discussion}

This deployment does not identify conditions under which AI helps governments. A single system, built one way, cannot show that any of its prior design choices was necessary. The four statements below are the principles the prototype was built to illustrate. They are hypotheses. Confirming or rejecting them requires evaluation that this internship has not done.

\paragraph{H1 (scope).}
An assistant is more likely to be usable in climate and energy briefing if it is limited to retrieval, comparison, and summarization over verified sources, and is forbidden from inventing measures or speaking in the Party's name. This restates the product brief. It is consistent with literature that treats traceability as central to whether LLM outputs are usable for policy decisions [2], but consistency with prior writing is not evidence from this cockpit.

\paragraph{H2 (institutional encoding).}
The same assistant is more likely to be usable if BAU-relative targets, conditionality, and district non-scoring are encoded the way ministries already understand them, rather than as absolute cuts from a base year. The transport example in Section~\ref{sec:setting} is a design rationale, not a measured reduction in misreading.

\paragraph{H3 (checkable numbers).}
Numeric claims should resolve to URLs a desk officer can open, and a near-miss fabrication should not receive those URLs. Section~\ref{sec:ndc-ai} specifies overwrite-or-flag instead of certify-by-tolerance. Hypothesis: an adversarial set of close-but-wrong queries will produce a lower false-verification rate under exact matching than under the $2.5\%$ rule. That rate has not been measured.

\paragraph{H4 (honesty partition).}
Live observations, official pledges, snapshot documents, and indicative screens should not be interchangeable, so a carbon-price slider cannot contaminate an inventory total. Table~\ref{tab:honesty} is the partition we implemented. We have not shown that users respect it, nor that a weaker partition would fail.

\paragraph{What remains unwarranted.}
We would be cautious about applications that generate policy recommendations or public communications without an explicit, checkable source trail, given the diplomatic and legal weight these documents carry. That caution is prior belief plus the known failure mode in Section~\ref{sec:ndc-ai}, not a comparative result. In high-stakes settings, an interface that refuses a question may be more valuable than one that always answers; that too is a hypothesis. Climate TRACE and national inventories remain incommensurable by construction. Indicative modules (finance, KCI analogies, 2030 trends) are useful for screening only if they are not treated as audited inputs to a submission. The generative-agent literature on public acceptability is a reminder that LLMs remain imperfect approximations of human response and should not substitute for participatory policy processes [4].

\section{Limitations and future work}
\label{sec:limitations}

This is a single deployed prototype, not a controlled study. We do not have systematic user studies of negotiator trust, query accuracy, or time savings relative to existing manual workflows. A June 2026 demonstration and executive recognition do not fill that gap. Immediate work is: (i)~replace \eqref{eq:progress-2030} with \eqref{eq:progress-t} in the shared progress engine and lock the new definition in tests; (ii)~implement overwrite-or-flag citation resolution and score it on an adversarial close-but-wrong query set, reporting false-verification and missed-citation rates; (iii)~compare assistant answers against expert-verified ground truth, including off-ledger questions. The document corpus is a dated export; passage coverage is incomplete relative to a Party's full policy stack. Extending live wiring beyond Uganda requires per-country NDC encodings, GADM mappings, and a repeated honesty audit---not only a new flag on a map. Linear BAU interpolation may not match any Party's official projection model. Finally, this is a planning and briefing tool. It is not software for preparing a Biennial Transparency Report or for replacing a national inventory agency.

\section{Conclusion}
\label{sec:conclusion}

The NDC Data Explorer is one scoped attempt at a citation-grounded cockpit on verified national submissions and independent emissions data, built with institutional stakeholders. The technical claim is modest and still untested at scale. A typed fact ledger plus exact citation handling will not make a language model a negotiator. The hypothesis is that it can make the model refuse to sound like one when the evidence is missing, provided progress is not inflated by a 2030 BAU comparator and provided near-miss numerals are not dressed with real sources. Whether that pattern helps governments, and whether the four design principles are necessary for that help, remains an open empirical question.

\begin{ack}
The author thanks Joaquim Barris, Data Science Manager at UNFCCC, for supervision, technical mentorship, and allowing this project to proceed; and Nicolas Henriksson, Policy Consultant at UNFCCC, for keeping the policy grounding sharp. Gavin McCormick, Executive Director of Climate TRACE and WattTime, and Peter Thomas, Senior Data Engineer at Climate TRACE, together with the teams at Climate TRACE and Planet, provided the emissions data that powers the cockpit. Kyra Appleby and Anne Jelmar Sietsma at Climate Policy Radar provided policy documentation and insight for the document corpus. Professor Kyle Whyte (School for Environment and Sustainability, University of Michigan) reviewed a draft of this manuscript. Any remaining errors are the author's.
\end{ack}

\section*{References}

\noindent [1] P. Lewis, E. Perez, A. Piktus, F. Petroni, V. Karpukhin, N. Goyal, H. K\"uttler, M. Lewis, W. Yih, T. Rockt\"aschel, S. Riedel, and D. Kiela, ``Retrieval-augmented generation for knowledge-intensive NLP tasks,'' in {\it Advances in Neural Information Processing Systems}, vol. 33, 2020, pp. 9459--9474.

\noindent [2] F. Larosa et al., ``Large language models in climate and sustainability policy: limits and opportunities,'' {\it Environmental Research Letters}, vol. 20, no. 7, p. 074032, 2025, arXiv:2502.02191.

\noindent [3] E. Nabavi, H. R. Maier, S. Razavi, A. Hindes, M. Howden, W. Grant, and S. Raman, ``On using large language models to support social research for climate action,'' {\it npj Climate Action}, vol. 5, no. 1, 2026.

\noindent [4] A. Manivannan, V. Spaiser, T. J. B. Cann, J. Evans, J. P. Everall, et al., ``Generative AI for climate governance and acceptability-constrained policy design,'' {\it npj Climate Action}, vol. 5, no. 1, p. 37, 2026.

\noindent [5] United Nations Framework Convention on Climate Change, {\it Paris Agreement}, United Nations, 2015, Article 4.

\noindent [6] Republic of Uganda, Ministry of Water and Environment, ``Updated nationally determined contribution (NDC),'' September 2022. Available: \url{https://unfccc.int/NDCREG}

\noindent [7] C. Gordon et al., ``Closing gaps in emissions monitoring with Climate TRACE,'' arXiv:2511.19277, 2025.

\noindent [8] Climate TRACE Coalition, ``Climate TRACE emissions inventory,'' Zenodo, 2026. Available: \url{https://doi.org/10.5281/zenodo.20414751}

\noindent [9] Climate Policy Radar, ``All document text data,'' Hugging Face dataset, 2026.

\noindent [10] Food and Agriculture Organization of the United Nations, ``NDC tracking tool, module 3: progress toward BAU-relative targets,'' 2024.

\end{document}
